%\clearpage
%\newgeometry{margin=0pt} % Apply margin only for this page
%\thispagestyle{empty}
%\begin{center}
%\includegraphics[width=0.9\textwidth, height=\paperheight, keepaspectratio]{./images/ganapa.jpg}
%\end{center}
%\restoregeometry % Restore original geometry settings
%\newpage

\clearpage
\newgeometry{margin=0pt} % Apply margin only for this page
\thispagestyle{empty}
%\begin{figure}
%\centering
%\includegraphics[width=0.9\textwidth, height=\paperheight, keepaspectratio]{./images/ganapa.jpg}
%\includegraphics[width=\paperwidth, height=\paperheight]{./images/ganapa.jpg}
%\end{figure}
%\restoregeometry % Restore original geometry settings
%\newpage
\newgeometry{top=8mm, bottom=5mm, left=10mm, right=10mm}
\thispagestyle{empty}
\begin{figure}[h!]
    \centering
    \begin{overpic}[width=0.8\paperwidth, height=0.9\paperheight]{../images/001.jpg}
        \put(8, 92){\color{white}\devfont गजाननं भूतगणादि सेवितं कपित्थजम्बूफलसार भक्षितम् ।
        }
        \put(8, 89){\color{white}\devfont उमासुतं शोकविनाशकारणं नमामि विघ्नेश्वर पादपङ्कजम् ॥}
        \put(12, 86){\color{white}\notosansdev\footnotesize gajānanaṁ bhūta-gaṇādi-sevitaṁ}
        \put(14, 84){\color{white}\notosansdev\footnotesize kapittha-jambū-phala-sāra bhakṣitam |}
        \put(12, 82){\color{white}\notosansdev\footnotesize umā-sutaṁ śoka-vināśa-kāraṇaṁ}
        \put(14, 80){\color{white}\notosansdev\footnotesize namāmi vighneśvara pāda-paṅkajam ‖}
      \end{overpic}
    %\caption{This is the standard figure caption below the image.}
    \label{fig:example}
\end{figure}

\restoregeometry
